% Options for packages loaded elsewhere
\PassOptionsToPackage{unicode}{hyperref}
\PassOptionsToPackage{hyphens}{url}
\documentclass[
  11pt,
]{article}
\usepackage{xcolor}
\usepackage[margin=1in]{geometry}
\usepackage{amsmath,amssymb}
\setcounter{secnumdepth}{-\maxdimen} % remove section numbering
\usepackage{iftex}
\ifPDFTeX
  \usepackage[T1]{fontenc}
  \usepackage[utf8]{inputenc}
  \usepackage{textcomp} % provide euro and other symbols
\else % if luatex or xetex
  \usepackage{unicode-math} % this also loads fontspec
  \defaultfontfeatures{Scale=MatchLowercase}
  \defaultfontfeatures[\rmfamily]{Ligatures=TeX,Scale=1}
\fi
\usepackage{lmodern}
\ifPDFTeX\else
  % xetex/luatex font selection
\fi
% Use upquote if available, for straight quotes in verbatim environments
\IfFileExists{upquote.sty}{\usepackage{upquote}}{}
\IfFileExists{microtype.sty}{% use microtype if available
  \usepackage[]{microtype}
  \UseMicrotypeSet[protrusion]{basicmath} % disable protrusion for tt fonts
}{}
\makeatletter
\@ifundefined{KOMAClassName}{% if non-KOMA class
  \IfFileExists{parskip.sty}{%
    \usepackage{parskip}
  }{% else
    \setlength{\parindent}{0pt}
    \setlength{\parskip}{6pt plus 2pt minus 1pt}}
}{% if KOMA class
  \KOMAoptions{parskip=half}}
\makeatother
\setlength{\emergencystretch}{3em} % prevent overfull lines
\providecommand{\tightlist}{%
  \setlength{\itemsep}{0pt}\setlength{\parskip}{0pt}}
\usepackage{bookmark}
\IfFileExists{xurl.sty}{\usepackage{xurl}}{} % add URL line breaks if available
\urlstyle{same}
\hypersetup{
  pdftitle={Action-Class Authority: Why the Impossibility of Complete Guardrails Forces a Reversibility-Graded Verification Layer for Agentic AI},
  pdfauthor={Mayur Agnihotri},
  hidelinks,
  pdfcreator={LaTeX via pandoc}}

\title{Action-Class Authority: Why the Impossibility of Complete
Guardrails Forces a Reversibility-Graded Verification Layer for Agentic
AI}
\author{Mayur Agnihotri}
\date{}

\begin{document}
\maketitle
\begin{abstract}
Production agentic AI systems fail at one architectural seam: a
prompt-injected, memory-poisoned, or supply-chain-compromised agent
retains valid credentials and reasons over its own action authority.
This paper argues that the response to that failure is not a stronger
guardrail but a differently-placed one. A formal impossibility result
establishes that no finite guardrail can be universally robust against
adversarial input, and four independent empirical studies from May 2026
measure agent-reasoning compromise at rates high enough that any
reasoning-based defense is statistically dominated. Read together, the
formal and empirical results force a specific architecture. Because
prevention cannot be completed and detection carries latency, the
actions safe to leave under detection alone are exactly the reversible
ones, while irreversible actions require authority established before
execution by a mechanism the agent cannot rewrite. We call this pattern
action-class authority: a verification layer that decouples a trusted,
design-time reversibility classification from agent runtime reasoning,
enforces it at a deterministic gate, and governs multi-step chains by
the worst-case reversibility class present anywhere in the chain. We
document the OWASP AISVS v1.0 verification requirements that
operationalize this pattern (C9.2.3, C9.2.4, C9.2.10), analyze three
canonical 2026 incidents, position the pattern against the 2026
literature on agent runtime enforcement and containment, and close with
open research questions.
\end{abstract}

\subsection{1. Introduction}\label{introduction}

The agentic AI deployment year of 2026 has crystallized a question the
field was able to defer through 2024 and 2025: where does the security
boundary for an autonomous agent live? Through the early agentic era,
the boundary was treated implicitly as a property of the model. Better
instruction-following, better safety training, better refusal
calibration, the argument went, would produce agents whose reasoning
could be trusted to refuse the actions an organization did not want to
authorize. Frontier model releases were evaluated on benchmarks that
measured the model's ability to recognize unsafe requests and decline
them. Agent deployments inherited the same assumption: if the model is
safe, the agent built on top is safe.

Two independent lines of evidence closed that argument in 2026, one
formal and one empirical.

The formal line removes the last reading of this as an engineering gap
that better models eventually close. Extending Goedel's incompleteness
and Chaitin's uncomputability results to AI safety, recent work shows
that no finite set of guardrails can be universally robust against
adversarial input: natural language is unbounded while any rule set is
finite, so complete coverage was never a target more effort reaches. It
is mathematically foreclosed {[}6{]}. The standard operational reading
of that result, instrument for continuous runtime detection, states the
consequence only halfway. If prevention cannot be completed, the
guardrail is bypassed eventually, and detection is an after-the-fact
signal with nonzero latency. For an action that can be undone,
detect-and-recover still closes the loop. For an action that cannot,
detection that fires a moment too late is indistinguishable from no
detection at all. So the impossibility result forces a step past moving
the boundary to the action layer. It partitions that layer by
reversibility: reversible actions can rest on detection, while
irreversible ones require authority established before execution by a
mechanism that does not depend on the guardrail holding.

The empirical line supplies the rates. Four independent publications in
May 2026, spanning academic systems security, industry framework
engineering, empirical production-scale evaluation, and academic
adversarial machine learning, reached the same architectural conclusion
through four methodologies: agent reasoning cannot be the safety
boundary. The measured compromise rates (95 to 99 percent injection at
the memory layer, 100 percent trust on poisoned knowledge graphs across
nine frontier models, 60 to 89 percent attacker-intended action rates
after retrieval) are what the impossibility looks like in production.
The boundary has to live at the action layer, and the impossibility
result tells us how that layer must be organized: by reversibility.

\subsubsection{1.1 Contributions}\label{contributions}

This paper makes four contributions:

\begin{enumerate}
\def\labelenumi{\arabic{enumi}.}
\item
  \textbf{The forcing argument.} We connect the formal impossibility of
  complete guardrails to a specific architectural consequence:
  reversibility-graded authority is structurally required, not
  defense-in-depth prudence. When prevention is provably incomplete and
  detection carries latency, irreversible actions cannot be made safe by
  detection and therefore require pre-execution authority independent of
  the guardrail. This derivation, and not the reversibility taxonomy
  itself, is the paper's primary claim.
\item
  \textbf{The verification pattern, spec-anchored.} We specify
  action-class authority as a verification-layer pattern and document
  the OWASP AISVS v1.0 controls that operationalize it: C9.2.3 (trusted
  reversibility classification), C9.2.4 (runtime enforcement by class),
  and C9.2.10 (worst-case reversibility class across a multi-step or
  multi-agent chain), with a worked threat model showing why the same
  agent cannot both propose and classify its action class.
\item
  \textbf{Worked case analysis.} We analyze three canonical 2026
  incidents (Codex Docker bind-mount, Meta Instagram AI Support
  Assistant, Red Hat npm Miasma supply chain) showing the architectural
  pattern is identical across coding agent, customer-facing assistant,
  and CI/CD pipeline.
\item
  \textbf{Positioning against the 2026 literature.} We situate
  action-class authority within the 2026 body of work on agent runtime
  enforcement, reversibility-tiered autonomy, and containment, and
  identify what that work leaves open: the derivation of
  reversibility-partitioning from impossibility, and the requirement
  that pre-execution authority be captured at an evidence vantage the
  agent cannot forge or omit, with binding strength scaled to
  irreversibility.
\end{enumerate}

The remainder of the paper is organized as follows. Section 2 states the
forcing argument. Section 3 describes the May 2026 empirical
convergence. Section 4 specifies the verification pattern and its AISVS
anchor. Section 5 presents three canonical cases. Section 6 compares
against LLM-guardian alternatives. Section 7 positions the pattern
against related work. Section 8 enumerates open research questions.
Section 9 concludes.

\begin{center}\rule{0.5\linewidth}{0.5pt}\end{center}

\subsection{2. The Forcing Argument}\label{the-forcing-argument}

The claim of this section is narrow and it is the paper's spine: given
the impossibility of complete guardrails and the latency of detection, a
reversibility partition of the action space is forced.

Let a guardrail be any mechanism that inspects a proposed action or its
context and decides whether to permit it. The impossibility result
{[}6{]} establishes that for any finite guardrail there exists
adversarial input that defeats it. This is not a statement about current
model quality; it is a statement about finite rule systems facing
unbounded input, in the lineage of Goedel and Chaitin. The empirical
convergence of Section 3 is the same statement measured in production.

Take the impossibility as given. Then, over a long enough horizon, the
guardrail is bypassed. The system's residual safety is whatever the
agent can reach in the window between bypass and response. Two facts
define that window. First, detection is an occurrence signal: it reports
that something happened, after it happened, with nonzero latency.
Second, actions differ in whether their effects can be withdrawn once
they occur.

Cross these two facts. For an action whose effects can be withdrawn,
detection with latency is sufficient: the loop closes by
detect-and-recover, and the residual exposure is bounded by recovery
time. For an action whose effects cannot be withdrawn, detection with
latency is not sufficient by construction: by the time the signal
arrives, the irreversible effect has committed, and detection that fires
a moment too late is indistinguishable from no detection. No improvement
in detection quality changes this, because the limiting quantity is
latency against irreversibility, not accuracy.

The consequence is a partition, not a preference. Reversible actions can
be governed by detection. Irreversible actions must be governed by
authority established before execution, and that authority cannot depend
on the same guardrail the impossibility result says will be bypassed. It
must be carried by a mechanism outside the agent's reasoning: a
classification the agent does not supply at runtime, evaluated by a gate
the agent does not control. This is the architecture the rest of the
paper specifies. The reversibility taxonomy is old; the forcing of it
from impossibility is the contribution.

A second consequence follows for evidence. If pre-execution authority is
the control for irreversible actions, the authorization must be captured
where the agent cannot forge or omit it, and the required strength of
that binding scales with how irreversible the action is. An
authorization asserted by the same component that can be compromised is
not authorization. We return to this in Section 4.3 and Section 7.

\begin{center}\rule{0.5\linewidth}{0.5pt}\end{center}

\subsection{3. The May 2026 Empirical
Convergence}\label{the-may-2026-empirical-convergence}

Four publications landed in May 2026 that, read together, settle the
agent-reasoning-as-boundary question through four independent
methodologies.

\subsubsection{3.1 Systems-Security
Framing}\label{systems-security-framing}

Christodorescu and Fernandes and collaborators argue from
systems-security first principles that the AI model must be treated as
an untrusted component, with security invariants enforced at the system
level rather than within the model's reasoning {[}1{]}. The argument is
structural: a model that can be prompted, poisoned, or contextually
shifted into producing arbitrary outputs cannot also be the mechanism
that decides which outputs are safe to act on.

\subsubsection{3.2 Operational Framework}\label{operational-framework}

Anthropic published a comprehensive operational framework for Zero Trust
in agentic AI systems {[}2{]}. The framework identifies the boundary
problem as a top-of-stack concern and introduces the
impossible-versus-tedious design test: controls whose value is friction
degrade against an adversary with unlimited patience, while controls
that remove a capability survive. The present work extends that
principle from the credential layer {[}11{]} to the action layer.

\subsubsection{3.3 Oracle Poisoning
Measurement}\label{oracle-poisoning-measurement}

Kereopa-Yorke and collaborators demonstrated that every tested
state-of-the-art model trusts moderately-sophisticated poisoned
knowledge-graph data at 100 percent {[}3{]}. The setup uses a
42-million-node code knowledge graph and nine state-of-the-art models.
Three attacker sophistication levels show that the
action-on-retrieved-content rate scales with sophistication while
detection rates do not.

\subsubsection{3.4 Sleeper Memory
Poisoning}\label{sleeper-memory-poisoning}

Pulipaka and collaborators reached the same conclusion at the memory
layer: 95 to 99 percent injection rates, 90 to 95 percent retrieval
rates on goal-adjacent queries, and 60 to 89 percent attacker-intended
action rates after retrieval {[}4{]}. A poisoned memory persists across
sessions; the agent's reasoning never sees the injection point because
the injection happened at write time, not read time.

\subsubsection{3.5 The Convergent
Conclusion}\label{the-convergent-conclusion}

Read together, the four constitute a structural argument: agent
reasoning is not a safety boundary because the reasoning component is
the same one an adversary can corrupt, and corruption rates at memory
and retrieval are high enough that any single-layer defense based on
reasoning is statistically dominated. This is the empirical face of the
forcing argument in Section 2.

\begin{center}\rule{0.5\linewidth}{0.5pt}\end{center}

\subsection{4. The Verification Pattern}\label{the-verification-pattern}

The pattern that operationalizes the forcing argument has three
requirements, all shipped in OWASP AISVS v1.0 (released 2026-06-24, C9
Orchestration and Agentic Action chapter) {[}5{]}.

\subsubsection{4.1 C9.2.3 Trusted Reversibility
Classification}\label{c9.2.3-trusted-reversibility-classification}

C9.2.3 (Level 2) requires that each high-impact action have a trusted
reversibility classification: read-only, reversible, externally
reversible, or irreversible. Trusted means the classification is not
supplied by the agent at runtime. The mechanism this paper specifies for
making the classification trusted is to declare it in the tool or action
manifest at design time, read by the gate rather than derived from the
agent's output. The four classes:

\begin{itemize}
\tightlist
\item
  \textbf{Read-Only (RO):} observes state without modifying it.
\item
  \textbf{Reversible (R):} modifies state within the system; recoverable
  via undo log, transaction abort, or compensating action.
\item
  \textbf{External-Reversible (ER):} affects external systems but is
  recoverable through identity verification or an out-of-band channel.
\item
  \textbf{Irreversible (I):} commits to durable state the user cannot
  recover from without out-of-band action; emails sent, funds
  transferred, packages published to public registries, account
  identifiers linked.
\end{itemize}

The classification lives in the manifest. The gate evaluates the
manifest. The agent reasons elsewhere. This is the architectural
separation the forcing argument requires.

\subsubsection{4.2 C9.2.4 Runtime Enforcement and C9.2.10 Worst-Case
Chain
Rule}\label{c9.2.4-runtime-enforcement-and-c9.2.10-worst-case-chain-rule}

C9.2.4 (Level 2) requires that the runtime enforce the classification by
blocking, requiring approval, or restricting actions based on impact and
reversibility. C9.2.10 (Level 3) requires that approval gates for
multi-step or multi-agent action chains enforce the highest-impact
reversibility class present anywhere in the chain. The chain rule
prevents the composition failure where individually low-gate steps chain
into an irreversible end state: read user email (RO), identify financial
accounts (RO), initiate password reset (ER), modify account credentials
(I) is a chain whose worst-case class is irreversible from the moment
the first step executes, not at the irreversible step.

\subsubsection{4.3 Threat Model and Evidence
Vantage}\label{threat-model-and-evidence-vantage}

The threat C9.2.3 addresses: when the same agent both proposes and
classifies an action, a prompt-injected or context-corrupted agent can
mislabel an irreversible action as reversible, and a gate reading the
runtime classification will permit it. The manifest-declared
classification was set at design time and cannot be rewritten by the
agent at runtime.

The forcing argument's second consequence (Section 2) sharpens this. The
gate's decision is only as sound as the evidence it reads. For the
irreversible class, the authorization must be captured at a vantage the
agent cannot forge or omit, and the binding strength should scale with
irreversibility. AISVS v1.0 reflects this gradient in its
higher-assurance controls: C9.2.8 binds approvals cryptographically to
action parameters, requester identity, execution context, and a
single-use nonce, and C9.2.9 isolates the approval key material from the
agent runtime {[}5{]}. Reversibility class thus governs not only whether
a gate fires but how strong the evidence behind it must be.

\begin{center}\rule{0.5\linewidth}{0.5pt}\end{center}

\subsection{5. Three Canonical Cases}\label{three-canonical-cases}

Three real-world cases from late May and early June 2026 demonstrate the
pattern.

\subsubsection{5.1 Codex Docker Bind-Mount (May 30,
2026)}\label{codex-docker-bind-mount-may-30-2026}

OpenAI's Codex coding agent, on a Linux host where the user was in the
docker group, chained four legitimate operations: enumerate group
memberships, recognize docker as root-equivalent, construct a bind-mount
of the host /etc into a container with write access, and execute a write
that overwrote a live configuration file. Each operation was permitted
by the tool authority; the composite was irreversible at the system
level. Under C9.2.3, ``write to /etc on a live host'' carries an
irreversible classification in the manifest; the deterministic gate
routes the action to human ratification regardless of how the agent
reasoned about the workflow, and C9.2.10 gates the whole chain by that
worst-case class.

\subsubsection{5.2 Meta Instagram AI Support Assistant (May 31,
2026)}\label{meta-instagram-ai-support-assistant-may-31-2026}

Meta's AI Support Assistant on Instagram had tool authority to modify
account state, including linking new email addresses. Attackers found
the prompt pattern that triggered an email-link change without identity
verification. Exploitation persisted for approximately four months
across thousands of accounts. Email-link to a primary identifier is
irreversible at the user level, because password-reset control follows
the email. A manifest-declared irreversible classification, evaluated by
a deterministic gate, routes the change to out-of-band identity
verification regardless of the attacker's prompt.

\subsubsection{5.3 Red Hat npm Miasma Supply Chain (June 1,
2026)}\label{red-hat-npm-miasma-supply-chain-june-1-2026}

The redhat-cloud-services npm scope was compromised with 32 packages
backdoored. The vector was not credential theft: attackers pushed orphan
commits adding a workflow with id-token write permissions, then used
OIDC trusted publishing to push through the legitimate path {[}8{]}.
Trusted publishing, designed to eliminate long-lived tokens, becomes a
misleading trust signal when pipeline access is compromised. Under
C9.2.3, ``publish to a public npm scope'' carries an irreversible
classification; C9.2.10 gates any chain reaching the publish step as
irreversible from the chain start.

\subsubsection{5.4 The Pattern Across
Cases}\label{the-pattern-across-cases}

Across a coding agent, a customer-facing assistant, and a CI/CD
pipeline, from three vendors and three technology stacks, the
composition failure is the same: legitimate per-step authority sequenced
into an irreversible composite that no reasoning-layer check reliably
stops. The missing layer in each is a trusted action-class declaration
evaluated by a deterministic gate, governed by the worst-case class
across the chain.

\begin{center}\rule{0.5\linewidth}{0.5pt}\end{center}

\subsection{6. Comparison with LLM-Guardian
Approaches}\label{comparison-with-llm-guardian-approaches}

An alternative architecture uses a second LLM as a guardian that
evaluates the primary agent's proposed actions before they execute. The
guardian approach has empirical attraction, but inherits the threat
model action-class authority is built to escape: a guardian whose
reasoning can be context-shifted or prompt-injected is structurally
identical to the agent it guards. The deterministic gate evaluating a
manifest-declared classification does not have this property, because
the classification is set at design time and is not rewritten at runtime
by any reasoning component.

This is not a rejection of LLM-guardian approaches, and the released
standard makes the relationship explicit. AISVS C9.2.6 requires that
AI-augmented review of planned high-risk actions add to, and not
replace, the deterministic policy gate, and C9.2.7 requires that this
review be protected against manipulation and not be bypassable through
prompt injection {[}5{]}. The standard therefore codifies the layering
this paper argues: the deterministic gate carries the structural
verification load; AI-augmented review adds context-aware checks above
it. The architectural question is which layer carries the load, not
whether AI-based checks have value.

\begin{center}\rule{0.5\linewidth}{0.5pt}\end{center}

\subsection{7. Related Work}\label{related-work}

\textbf{Foundations.} The separation of a declared capability from the
actor that exercises it is the core of capability-based security
{[}10{]} and has parallels in transactional-database reversibility,
formal verification of reversible computation, and digital forensic
chain of custody. Action-class authority applies that separation to
agent action manifests.

\textbf{The impossibility spine.} The formal basis for Section 2 is
Vassilev's result that no finite guardrail is universally robust against
adversarial input, extending Goedel and Chaitin to AI safety {[}6{]}. A
related archival note reaches a similar conclusion on containment from
logical and computational constraints {[}22{]}. Moon and Varshney give a
complementary formal result, verifying safety over a typed action
boundary in a way that is invariant to model capability {[}14{]}. These
establish that prevention is bounded and that a fixed boundary can be
verified. None derives the reversibility partition of Section 2, in
which impossibility plus detection latency forces pre-execution
authority specifically for the irreversible class.

\textbf{Reversibility-tiered execution.} Two 2026 lines classify agent
actions by reversibility. Sahoo's Controllability Trap scores each tool
call on an irreversibility scale and enforces a cumulative
irreversibility budget before re-authorization {[}12{]}. Ramaswamy and
Miaosheng route operations through execution gears with utility-gated
dispatch and stability guarantees for cyber-physical systems {[}13{]}.
Both tier autonomy by reversibility as prudent design; neither derives
the tiering from a formal impossibility result, and both key the gate on
utility or a cumulative budget rather than a trusted, non-agent-supplied
class evaluated per action.

\textbf{The nearest neighbors.} Two papers each cover part of this
paper's synthesis and are the sharpest comparisons. Zhai, Li and Wang's
Revisable by Design formalizes the same four-way reversibility taxonomy
and proves that conflicting irreversible actions make full specification
satisfaction impossible, a property of the action space rather than the
algorithm {[}7{]}. That composition result is adjacent to the worst-case
chain rule of Section 4.2, and it differs on two axes. Its impossibility
concerns specification satisfaction and rollback cost in a
streaming-intervention model with mid-execution rollback, not the
guardrail-bypass impossibility this paper uses to force pre-execution
authority; and it does not treat the classification as trusted
non-agent-supplied metadata or address the evidence vantage of the
authority that gates an irreversible action. Dobrin and Chmiel's
Unfireable Safety Kernel supplies the other half: external pre-action
enforcement on a structurally sole path, fail-closed, with externalized
signed evidence verifiable outside the agent's trust boundary {[}9{]}.
This is the unforgeable evidence vantage Section 2's second consequence
requires. It differs in that its authority is uniform, applying
identical logic to every safety-critical operation with no grading by
reversibility, and it is derived from architectural properties rather
than an impossibility result, with no worst-case-across-chain rule. The
contribution of this paper is the composition these two leave apart:
guardrail-impossibility used to force a reversibility partition,
pre-execution authority at an unforgeable evidence vantage whose
strength scales to that partition, and multi-step chains governed by the
least-reversible step. Metere's Skills as Verifiable Artifacts gates a
human-in-the-loop check on a skill's verification level and observes
that gating every irreversible call degrades into rubber-stamping
{[}19{]}; scaling authority by a reversibility class rather than a
uniform irreversible flag is the response to exactly that degradation.

\textbf{Runtime enforcement and separation.} A further body of 2026 work
enforces action policies at runtime without a reversibility grading:
neuro-symbolic constraint satisfaction {[}15{]}, correct-by-construction
verified policies with per-action monitoring {[}16{]}, DSL-specified
runtime rules {[}17{]}, and hard cognitive-executive separation in which
thinking agents may not act {[}18{]}. These reduce unsafe execution but
gate on generic policy or a binary act/no-act split rather than a
trusted reversibility class scaled to how irreversible the action is,
and none governs multi-step chains by their worst-case class. Sah and
colleagues quantify a related horizon-dependent tradeoff between
verification cost and safety in tool-using agents {[}20{]}.

\textbf{Supply-chain trust inheritance.} The manifest-integrity
assumption underlying action-class authority connects to trust
inheritance across build and skill layers, where a declared claim is
trusted for its position in a registry rather than a binding the actor
cannot forge {[}21{]}.

\begin{center}\rule{0.5\linewidth}{0.5pt}\end{center}

\subsection{8. Open Research Questions}\label{open-research-questions}

\begin{enumerate}
\def\labelenumi{\arabic{enumi}.}
\item
  \textbf{Input-layer corruption vs.~action-layer gate.} When the
  manifest itself is corrupted at distribution time, what is the
  residual exposure of an otherwise-correct gate? Hypothesis: exposure
  is bounded by manifest integrity-binding (signature, chain-of-custody
  at install); measurement against signed vs.~unsigned manifest
  distributions would settle the threshold.
\item
  \textbf{Multi-step information-flow control.} The chain rule gates
  chains; it does not detect cross-chain information flows where
  reversibility-graded data moves between chains. Hypothesis:
  information-flow labeling with reversibility class as a label
  dimension extends the chain rule to cross-chain composition.
\item
  \textbf{Runtime drift detection.} When an agent's action distribution
  shifts toward higher-risk classes over time, is the drift detectable
  from action-class telemetry alone, or does it require correlated
  identity, posture, and content signals?
\item
  \textbf{Reversibility-classification disagreement.} When vendor
  manifest and enterprise policy classify the same action differently,
  what resolves the disagreement? Hypothesis: worst-case wins is the
  safe default; case studies of real conflicts would test it.
\item
  \textbf{Reversibility-class evolution.} When an originally-reversible
  action becomes irreversible due to downstream propagation, how is the
  classification updated, and what is the gate behavior in the
  transition window?
\end{enumerate}

\begin{center}\rule{0.5\linewidth}{0.5pt}\end{center}

\subsection{9. Conclusion}\label{conclusion}

The empirical evidence and the formal result point the same way.
Prevention cannot be completed, so the guardrail is bypassed eventually;
detection reports the bypass after it happens, with latency. For
reversible actions that is enough. For irreversible actions it is not,
by construction, and no improvement in detection closes the gap. The
defense that holds partitions the action space by reversibility and
moves the authority for irreversible actions ahead of execution, onto a
classification the agent does not supply and a gate the agent does not
control. The OWASP AISVS v1.0 controls C9.2.3, C9.2.4, and C9.2.10
operationalize this partition; C9.2.6 through C9.2.9 layer AI-augmented
review and cryptographic binding above it. The conversation is no longer
about whether agent reasoning is trustworthy enough to act on. The
conversation is about where the gate sits, what it evaluates against,
and how strong that evidence must be for the actions that cannot be
undone.

\begin{center}\rule{0.5\linewidth}{0.5pt}\end{center}

\subsection{References}\label{references}

{[}1{]} M. Christodorescu, E. Fernandes, A. Hooda, S. Jha, J. Rehberger,
K. Chaudhuri, X. Fu, K. Shams, et al., ``Agent Security is a Systems
Problem,'' arXiv:2605.18991, 2026.

{[}2{]} Anthropic, ``Zero Trust for AI Agents,'' enterprise framework,
May 2026.

{[}3{]} B. Kereopa-Yorke et al., ``Oracle Poisoning: Corrupting
Knowledge Graphs to Weaponise AI Agent Reasoning,'' arXiv:2605.09822,
2026.

{[}4{]} S. Pulipaka, S. Hlebik, L. Raghav, S. Abdelnabi, V. Raina, I.
Sheth, M. Fritz, ``Hidden in Memory: Sleeper Memory Poisoning in LLM
Agents,'' arXiv:2605.15338, 2026.

{[}5{]} OWASP AISVS Project, ``AI Security Verification Standard,'' v1.0
(released 2026-06-24), C9 Orchestration and Agentic Action, controls
C9.2.3, C9.2.4, C9.2.6, C9.2.7, C9.2.8, C9.2.9, C9.2.10.
github.com/OWASP/AISVS.

{[}6{]} A. Vassilev, ``Robust AI Security and Alignment: A Sisyphean
Endeavor?,'' IEEE Security \& Privacy, 2026 (arXiv:2512.10100).

{[}7{]} Z. Zhai, M. Li, X. Wang, ``Revisable by Design: A Theory of
Streaming LLM Agent Execution,'' arXiv:2604.23283, 2026.

{[}8{]} Aikido Security, ``Analysis of the redhat-cloud-services npm
Compromise (Miasma),'' June 1, 2026.

{[}9{]} S. Dobrin, L. Chmiel, ``The Unfireable Safety Kernel:
Execution-Time AI Alignment for AI Agents and Other Escapable AI
Systems,'' arXiv:2606.26057, 2026.

{[}10{]} M. S. Miller, ``Robust Composition: Towards a Unified Approach
to Access Control and Concurrency Control,'' Ph.D.~dissertation, Johns
Hopkins University, 2006.

{[}11{]} D. Hardt, ``Anthropic's Zero Trust for AI Agents Sets the Right
Test: The Bearer Token Fails It,'' Hello Blog, June 2026.

{[}12{]} S. Sahoo, ``The Controllability Trap: A Governance Framework
for Military AI Agents,'' arXiv:2603.03515, 2026 (ICLR 2026 Workshop on
Agents in the Wild).

{[}13{]} S. Ramaswamy, W. Miaosheng, ``Managed Autonomy at Runtime:
Gear-Based Safety and Governance for Single- and Multi-Agent
Cyber-Physical Systems,'' arXiv:2607.00334, 2026.

{[}14{]} R. Moon, L. R. Varshney, ``Containment Verification: AI Safety
Guarantees Independent of Alignment,'' arXiv:2605.09045, 2026.

{[}15{]} B. Wu, W. Zhang, K. Chen, H. Fang, N. Yu, ``Provably Secure
Agent Guardrail,'' arXiv:2605.29251, 2026.

{[}16{]} ``VeriGuard: Enhancing LLM Agent Safety via Verified Code
Generation,'' arXiv:2510.05156, 2025.

{[}17{]} H. Wang, C. M. Poskitt et al., ``AgentSpec: Customizable
Runtime Enforcement for Safe and Reliable LLM Agents,''
arXiv:2503.18666, 2025.

{[}18{]} J. Fokou, ``Parallax: Why AI Agents That Think Must Never
Act,'' arXiv:2604.12986, 2026.

{[}19{]} A. Metere, ``Skills as Verifiable Artifacts: A Trust Schema and
a Biconditional Correctness Criterion for Human-in-the-Loop Agent
Runtimes,'' arXiv:2605.00424, 2026.

{[}20{]} T. Sah, V. Srivastava, D. Sah, K. Jordan, ``The Verifier Tax:
Horizon Dependent Safety Success Tradeoffs in Tool Using LLM Agents,''
ACM CAIS 2026 (arXiv:2603.19328).

{[}21{]} M. Agnihotri, ``Trust Inheritance Across Layers: A Cross-Domain
Pattern Analysis of June 2026 Supply-Chain Compromises,'' SSRN preprint,
DOI 10.2139/ssrn.6931278, 2026.

{[}22{]} S. Haider, ``The Impossibility of AI Containment: Logical,
Mathematical, Computational and Philosophical Constraints,''
philsci-archive \#24223, 2026.

\end{document}
